\section*{الفصل \ref{ch:b1:electrostatics-gauss} --- الكهرباء الساكنة: المجال وقانون غاوس}

\begin{solution}{exo:b1:electrostatics-gauss:1}
$F_e = 8.99 \times 10^9 \times (1.6 \times 10^{-19})^2/(2 \times 10^{-15})^2 = \qty{58}{N}$؛
و $F_g = 6.67 \times 10^{-11} \times (1.67 \times 10^{-27})^2/(4 \times 10^{-30}) =
\qty{4.6e-35}{N}$؛ والنسبة $1.2 \times 10^{36}$. وتتماسك النوى لأن التآثر
القوي، على مدى الفمتومتر، يغلب حتى هذا التنافر.
\end{solution}

\begin{solution}{exo:b1:electrostatics-gauss:2}
$E = 8.99 \times 10^9 \times 10^{-6}/1 = \qty{9.0}{kV/m}$. والمجال المتجه نحو الأسفل هو
مجال شحنة سالبة: ومنه $Q = -4\pi\varepsilon_0R^2E = -1.11 \times 10^{-10} \times
(6.37 \times 10^6)^2 \times 100 = \qty{-4.5e5}{C}$.
\end{solution}

\begin{solution}{exo:b1:electrostatics-gauss:3}
على المحور، من أجل $|x| > a$: $E_x = \dfrac{q}{4\pi\varepsilon_0}\Bigl[\dfrac{1}{(x - a)^2} +
\dfrac{1}{(x + a)^2}\Bigr]$، مبتعدًا. وعلى المحور $y$ تتلاشى المركّبات على $x$:
فيكون $E_y = \dfrac{2qy}{4\pi\varepsilon_0(y^2 + a^2)^{3/2}}$؛ و $\dd E_y/\dd y \propto
(y^2 + a^2) - 3y^2 = 0$ عند $y = a/\sqrt2$. وينعدم المجال عند المبدأ؛
والخطوط التي تغادر كل شحنة نحو الأخرى تنحرف مبتعدةً على امتداد
$\pm y$ --- وهي نقطة سرج.
\end{solution}

\begin{solution}{exo:b1:electrostatics-gauss:4}
بين الشحنتين يتجه المجالان كلاهما نحو $-q$: فلا انعدام. وعلى يسار
$+2q$ يكون $2q/x^2 > q/(x + a)^2$ دائمًا. ووراء $-q$، على بعد $d$
منها: $2q/(a + d)^2 = q/d^2$، ومنه $a + d = \sqrt2\,d$ و $d = a/(\sqrt2 - 1) = 2.4a$.
ويغادر خطان $+2q$ مقابل كل خط ينتهي على $-q$؛ ويذهب النصف الآخر إلى
اللانهاية، حيث يبدو الزوج شحنةً $+q$.
\end{solution}

\begin{solution}{exo:b1:electrostatics-gauss:5}
حلقة نصف قطرها $a$ وعرضها $\dd a$: $\dd q = 2\pi\sigma a\,\dd a$ و $\dd E_z =
\dfrac{\sigma z}{2\varepsilon_0}\dfrac{a\,\dd a}{(a^2 + z^2)^{3/2}}$؛ وبالمكاملة:
$\dfrac{\sigma z}{2\varepsilon_0}\Bigl[-\dfrac{1}{\sqrt{a^2 + z^2}}\Bigr]_0^R = \dfrac{\sigma}{2\varepsilon_0}
\Bigl(1 - \dfrac{z}{\sqrt{z^2 + R^2}}\Bigr)$. وعند $z \ll R$: $\sigma/2\varepsilon_0$، أي
المستوي اللانهائي. وعند $z \gg R$: $z/\sqrt{z^2 + R^2} \approx 1 - R^2/2z^2$، ومنه $E_z
\approx \sigma R^2/4\varepsilon_0z^2 = Q/4\pi\varepsilon_0z^2$، أي الشحنة النقطية.
\end{solution}

\begin{solution}{exo:b1:electrostatics-gauss:6}
$Q = 79e = \qty{1.26e-17}{C}$؛ و $E(R) = 8.99 \times 10^9 \times 1.26 \times 10^{-17}/
(7 \times 10^{-15})^2 = \qty{2.3e21}{V/m}$؛ وعند $R/2$ النصف: \qty{1.2e21}{V/m}؛ وعند
\qty{1}{pm}: $1.14 \times 10^{-7}/10^{-24} = \qty{1.1e17}{V/m}$ --- أي $10^{15}$ ضعف
مجال انهيار الهواء، على مسافة تعيش عندها الإلكترونات.
\end{solution}

\begin{solution}{exo:b1:electrostatics-gauss:7}
أسطوانة نصف قطرها $r$ وطولها $h$: $2\pi rh\,E = \lambda h/\varepsilon_0$، ومنه $E =
\lambda/2\pi\varepsilon_0r$. ويقع التفريغ الهالي حين $\lambda = 2\pi\varepsilon_0 rE = 2\pi \times 8.85 \times
10^{-12} \times 0.01 \times 3 \times 10^6 = \qty{1.7}{\micro C/m}$؛ وعند \qty{5}{m} يكون
$E = 3 \times 10^6 \times 0.01/5 = \qty{6}{kV/m}$.
\end{solution}

\begin{solution}{exo:b1:electrostatics-gauss:8}
$E = \sigma/\varepsilon_0 = 10^{-6}/8.85 \times 10^{-12} = \qty{1.1e5}{V/m}$. ولا تؤثر صفيحة
بقوة صافية في نفسها (فإسهامات مجالها الخاص تتلاشى
مثنى مثنى)؛ ومجال الصفيحة الأخرى عند موضعها هو $\sigma/2\varepsilon_0$، ومنه
$F = Q\sigma/2\varepsilon_0 = 10^{-6} \times 10^{-6}/(2 \times 8.85 \times 10^{-12}) =
\qty{5.6e-2}{N}$، وهي جاذبة؛ والضغط $\sigma^2/2\varepsilon_0 = \qty{0.056}{Pa}$.
\end{solution}

\begin{solution}{exo:b1:electrostatics-gauss:9}
عند $r < a$: تستقر شحنة الناقل الداخلي على سطحه، فلا شحنة
في الداخل ويكون $E = 0$. وعند $a < r < b$: $E = \lambda/2\pi\varepsilon_0r$. وعند $r > b$: الشحنة
المحصورة $\lambda - \lambda = 0$، ومنه $E = 0$ --- فالكبل لا يشعّ أيّ مجال
ساكن. وعند $r = a$: $2 \times 8.99 \times 10^9 \times 10^{-8}/(5 \times 10^{-4}) =
\qty{3.6e5}{V/m}$.
\end{solution}

\begin{solution}{exo:b1:electrostatics-gauss:10}
$\rho = 3Q/4\pi(R_2^3 - R_1^3)$. وعند $r < R_1$: $E = 0$. وعند $R_1 < r < R_2$: $4\pi r^2E =
\rho\,\tfrac43\pi(r^3 - R_1^3)/\varepsilon_0$، ومنه $E = \dfrac{Q}{4\pi\varepsilon_0r^2}\,
\dfrac{r^3 - R_1^3}{R_2^3 - R_1^3}$. وعند $r > R_2$: $Q/4\pi\varepsilon_0r^2$. وعند $R_1$ يعطي
الجانبان $0$؛ وعند $R_2$ يعطيان $Q/4\pi\varepsilon_0R_2^2$. وفي القشرة يتزايد $E
\propto r - R_1^3/r^2$، فيكون الأعظمي عند $R_2$، ثم
يتناقص بنسبة $1/r^2$. وفي التجويف المجال معدوم: فلا تحسّ شحنة
هناك بأيّ قوة --- أينما استقرت.
\end{solution}

\begin{solution}{exo:b1:electrostatics-gauss:11}
(أ) اللبّ $\tfrac43\pi(3.48 \times 10^6)^3 \times 11000 = \qty{1.94e24}{kg}$؛ والوشاح
$\tfrac43\pi[(6.37 \times 10^6)^3 - (3.48 \times 10^6)^3] \times 4500 = \qty{4.08e24}{kg}$؛
والمجموع \qty{6.0e24}{kg}، بفارق دون $1\%$. (ب) وعند السطح $GM/R^2 = \qty{9.9}{m/s^2}$؛
وعند الحدّ $6.67 \times 10^{-11} \times 1.94 \times 10^{24}/(3.48 \times 10^6)^2 =
\qty{10.7}{m/s^2}$. (ج) ومع $g = GM(r)/r^2$ و $\dd M = 4\pi r^2\rho_m\,\dd r$:
$\dd g/\dd r = 4\pi G\rho_m - 2GM(r)/r^3 = 4\pi G[\rho_m - \tfrac23\bar\rho(r)]$،
حيث $\bar\rho(r) = M(r)/\tfrac43\pi r^3$ متوسط الكثافة تحت $r$.
وتتزايد $g$ مع العمق (أي $\dd g/\dd r < 0$) إذا وفقط إذا كان $\rho_m < \tfrac23\bar\rho$. وتحت
السطح مباشرةً يكون $\bar\rho = 5500$ و $\tfrac23\bar\rho = 3700 < 4500$: فتهبط $g$
هبوطًا طفيفًا أولًا؛ وعند حدّ اللبّ يكون $\bar\rho = 11000$ و $7300 >
4500$: فترتفع هناك إلى $10.7$.
\end{solution}

\begin{solution}{exo:b1:electrostatics-gauss:12}
(أ) في الداخل $E = er/4\pi\varepsilon_0R^3$ مبتعدًا؛ وعلى الإلكترون $\vect F =
-\dfrac{e^2}{4\pi\varepsilon_0R^3}\,\vect r$: أي نابض. (ب) و $\omega = \sqrt{e^2/4\pi\varepsilon_0mR^3}
= \sqrt{2.31 \times 10^{-28}/(9.11 \times 10^{-31} \times 10^{-30})} = \qty{1.6e16}{rad/s}$
و $f = \qty{2.5e15}{Hz}$ و $\lambda = c/f = \qty{120}{nm}$. (ج) وهو قريب على نحو مدهش
من الخط \qty{122}{nm}: فحجم الذرة وشدتها صحيحان.
لكن النموذج يعطي تواترًا واحدًا لا سلسلة، والنواة ليست
كرة قطرها \qty{0.1}{nm} بل نقطة أبعادها $10^{-15}$~متر --- وكان إلكترون كلاسيكي
يدور ليشعّ طاقته كلها في نانوثوانٍ. والعلاج
في \cref{ch:b1:quantum-introduction}.
\end{solution}

\begin{solution}{pb:b1:electrostatics-gauss:1}
\textbf{1.} $\vect F = -e\vect E$: أي نحو الأعلى. والمجال المتجه نحو الأسفل عند السطح
هو مجال شحنة سالبة: فالأرض سالبة.

\textbf{2.} كرة نصف قطرها $R_E$، مع $E_n = -100$: $Q = 4\pi\varepsilon_0R_E^2E_n =
1.11 \times 10^{-10} \times 4.06 \times 10^{13} \times (-100) = \qty{-4.5e5}{C}$.

\textbf{3.} $\sigma = \varepsilon_0E_n = \qty{-8.9e-10}{C/m^2}$؛ و $8.9 \times 10^{-10}/1.6 \times
10^{-19} = 5.5 \times 10^9$ إلكترونًا في المتر المربع، أي $5.5 \times 10^5$ في
السنتيمتر المربع.

\textbf{4.} $F = 100e \times 100 = \qty{1.6e-15}{N}$؛ والثقل $10^{-9} \times 9.8 =
\qty{9.8e-9}{N}$: أي أكبر ستة ملايين مرة.

\textbf{5.} النواظم الخارجية: في الأعلى $E_n = -5$، وفي الأسفل $E_n = +100$، وعلى الجوانب
صفر: ومنه $\Phi = \qty{95}{V.m}$، فتكون $Q = 95\varepsilon_0 = \qty{+8.4e-10}{C}$ في المتر
المربع --- أي شحنة فضائية موجبة تكاد تلاشي شحنة الأرض
$-8.9 \times 10^{-10}$.

\textbf{6.} $\rho = 8.4 \times 10^{-10}/10^4 = \qty{8.4e-14}{C/m^3}$: أي $5 \times 10^5$
\omterm{def:b1:electrostatics-gauss:charge}{شحنة أولية} في المتر المكعب؛ ومقابل $10^9$ أيون من كل إشارة،
يكون الاختلال $5 \times 10^{-4}$.

\textbf{7.} $J = \gamma E_0 = \qty{2e-12}{A/m^2}$؛ و $I = J \times 5.1 \times 10^{14} =
\qty{1.0}{kA}$؛ و $Q/I = 4.5 \times 10^5/1000 = \qty{450}{s}$. فكانت الأرض لتصير
متعادلة في ثماني دقائق: فثمة ما يعيد شحنها باستمرار.

\textbf{8.} يقع المحور في كل مستوٍ يحويه، وكلها مستويات
تناظر: فيكون $\vect E$ محوريًّا. ويعطي كل $\dd q$ المقدار $\dd q/4\pi\varepsilon_0(a^2 + z^2)$
في جهته هو، وجزؤه المحوري هو النسبة
$z/\sqrt{a^2 + z^2}$: ومنه $E_z = qz/4\pi\varepsilon_0(a^2 + z^2)^{3/2}$.

\textbf{9.} $\dd q = 2\pi\sigma a\,\dd a$: ومنه $E_z = \dfrac{\sigma z}{2\varepsilon_0}\displaystyle\int_0^R
\dfrac{a\,\dd a}{(a^2 + z^2)^{3/2}} = \dfrac{\sigma}{2\varepsilon_0}\Bigl(1 - \dfrac{z}{\sqrt{z^2
+ R^2}}\Bigr)$.

\textbf{10.} عند $z \ll R$: $\sigma/2\varepsilon_0$، أي المستوي اللانهائي (وهي قيمة غاوس)؛
وعند $z \gg R$: $\sigma R^2/4\varepsilon_0z^2 = Q/4\pi\varepsilon_0z^2$، أي الشحنة النقطية.

\textbf{11.} $\sigma = -40/(\pi \times 9 \times 10^6) = \qty{-1.4e-6}{C/m^2}$؛ و $\sigma/2\varepsilon_0 =
\qty{8.0e4}{V/m}$؛ وعند $z = \qty{3}{km}$ يكون $1 - 3/\sqrt{18} = 0.29$: ومنه $E =
\qty{2.3e4}{V/m}$، متجهًا نحو \emph{الأعلى} إلى الشحنة السالبة.

\textbf{12.} القرص العلوي، عند $z = \qty{9}{km}$: $8.0 \times 10^4 \times (1 - 9/\sqrt{90}) =
\qty{4.1e3}{V/m}$ نحو الأسفل. والصافي $\qty{1.9e4}{V/m}$ نحو الأعلى، ونحو $4 \times 10^4$
مع شحن الأرض المستحثة: وهي رتبة \qty{10}{kV/m}
المقيسة الصحيحة (فالشحنة الحقيقية أكثر انتشارًا ومحجوبة جزئيًّا)؛
وهو معكوس مجال الطقس الصافي وأكبر منه مئة مرة.

\textbf{13.} السحابة كلها: شحنتها الصافية معدومة، \omterm{def:b1:electrostatics-gauss:flux}{فالتدفق} معدوم. والقرص السفلي وحده:
$-Q/\varepsilon_0 = -40/8.85 \times 10^{-12} = \qty{-4.5e12}{V.m}$.

\textbf{14.} يبعد كل قرص \qty{3}{km} ويعطي $2.3 \times 10^4$؛ والعلوي
($+$) يدفع نحو الأسفل، والسفلي ($-$) يجذب نحو الأسفل: فيكون \qty{4.7e4}{V/m}
نحو الأسفل --- أي دون الانهيار بثلاثين مرة.

\textbf{15.} $Q/4\pi\varepsilon_0R^2 = 8.99 \times 10^9 \times 40/(2.5 \times 10^5) =
\qty{1.4e6}{V/m}$: فلا يبلغ الانهيارَ إلا شحنة محشورة في نصف كيلومتر،
والمجالات المقيسة داخل السحب أضعف عشر مرات. والمرشحان هما
التعزيز الموضعي عند القطرات والجليد، أو الإلكترونات الجامحة التي تبذرها الأشعة
الكونية --- فكيف يبدأ الوميض ما زال موضع جدال.

\textbf{16.} $I = 5/(5 \times 10^{-5}) = \qty{1e5}{A}$.

\textbf{17.} $W = QEd = 5 \times 10^5 \times 3000 = \qty{1.5e9}{J}$ (أي \qty{420}{kW h})؛
و $P = 1.5 \times 10^9/5 \times 10^{-5} = \qty{3e13}{W}$ --- أي مرة ونصف
استطاعة البشرية كلها، لخمسين ميكروثانية.

\textbf{18.} الحجم $\pi(0.02)^2 \times 3000 = \qty{3.8}{m^3}$، والكتلة \qty{4.5}{kg}؛
و $\Delta T = 1.5 \times 10^9/(4.5 \times 720) = \qty{4.6e5}{K}$ لو لم يفلت شيء.
لكن القناة تنفجر نحو الخارج (والصدمة هي الرعد)، وتشعّ ضوءًا
وموجات راديوية، وتؤيّن الهواء وتفكّكه: فتسخين القناة النهائية
إلى $3 \times 10^4$~\omterm{def:b1:kinetic-theory:temperature}{كلفن} جزء صغير من الفاتورة.

\textbf{19.} \qty{2.9}{s} لكل كيلومتر. ويغادر الصوت كل نقطة من قناة
طولها \qty{3}{km} في آنٍ واحد، من مسافات تختلف بالكيلومترات: فيصل
موزّعًا على ثوانٍ عدة --- وهو الهدير.

\textbf{20.} $50 \times 5 = \qty{250}{A}$ تُحمَل إلى الأرض؛ ومع تيارات
العواصف الأهدأ (التفريغ من الأطراف تحت السحب، والمطر المشحون)
يبلغ المجموع كيلوأمبير الجزء~I: فالعواصف مولّدات دارة
عالمية طريقُ عودتها هواءُ الطقس الصافي.

\textbf{21.} $\Phi_{\text{الخارج}}(\vect g) = -4\pi GM_{\text{الداخل}}$: فمن أجل كتلة نقطية
يكون $g = -GM/r^2$ على كرة مساحتها $4\pi r^2$، \omterm{def:b1:electrostatics-gauss:flux}{والتدفق} $-4\pi GM$، مستقلًّا
عن $r$؛ والباقي يتبع كما في \omterm{thm:b1:electrostatics-gauss:coulomb}{قانون كولوم}.

\textbf{22.} $g(r) = g_0r/R_E$؛ وعند $R_E/2$: \qty{4.9}{m/s^2}.

\textbf{23.} $M_{\text{اللبّ}} = \qty{1.94e24}{kg}$، ومنه $g = GM_{\text{اللبّ}}/R_c^2 =
\qty{10.7}{m/s^2}$ --- أي أكبر مما عند السطح: فالوشاح فوقه
أخفّ من متوسط ما تحته.

\textbf{24.} $m\ddot x = -mg_0x/R_E$: ومنه $\omega = \sqrt{g_0/R_E} = \qty{1.24e-3}{rad/s}$
و $T = 2\pi/\omega = \qty{5060}{s} = \qty{84}{min}$ و $v_{\max} = \omega R_E =
\qty{7.9}{km/s}$ --- وهي سرعة مدار قمر صناعي ملامس، ودوره
هو هو.

\textbf{25.} حلّ محلّ المجموع على كل شحنة من شحن الأرض (في I)،
وعلى أيونات الهواء (في I)، وعلى كل قشرة صخرية (في IV) ---
وفي~II عيّن حدود المجموع على القرص. واحفظ: \qty{-5e5}{C} و
\qty{1}{kA} للأرض في الطقس الصافي؛ و\qty{40}{C} و\qty{e4}{V/m} تحت
العاصفة؛ و\qty{e5}{A} و\qty{e9}{J} في الضربة؛ و\qty{10.7}{m/s^2} عند حدّ
اللبّ.
\end{solution}
